\section{Methodology: Standard-Deviation and RMS Scaling}
\label{submission:sec:method}

We now describe our proposed merging frameworks: \textbf{Standard-Deviation Scaling (SD-Scale)} and its mathematically superior, stable variant \textbf{Root-Mean-Square Scaling (RMS-Scale)}. We begin with the standard-deviation formulation, identify its mathematical vulnerabilities, present the RMS-Scale solution, and conclude with complexity and implementation details.

\subsection{Formulation and Notation}
Let $\Theta_{\text{pre}} \in \mathbb{R}^N$ be the parameters of a shared pretrained base model, and let $\{\Theta_1, \dots, \Theta_K\}$ be the parameters of $K$ independently fine-tuned expert models, where each $\Theta_k \in \mathbb{R}^N$ is initialized from $\Theta_{\text{pre}}$. Let $\Theta_{\text{pre}}^l$ and $\Theta_k^l$ denote the weight or bias tensor at a specific layer $l \in \{1, \dots, L\}$, where $L$ is the total number of parameter layers in the network. The number of parameters in layer $l$ is denoted by $N^l$.

For each task $k$ and layer $l$, the parameter update from pretraining is represented by the task vector:
\begin{equation}
\tau_k^l = \Theta_k^l - \Theta_{\text{pre}}^l
\end{equation}

\subsection{Standard-Deviation Scaling (SD-Scale)}
Standard Task Arithmetic~\cite{ilharco2022editing} merges models by simply averaging task vectors: $\tau_{\text{merged}}^l = \frac{1}{K}\sum_{k=1}^K \tau_k^l$. However, if task $A$ undergoes highly active updates during fine-tuning while task $B$ undergoes small, fine-grained adjustments, then $\text{std}(\tau_A^l) \gg \text{std}(\tau_B^l)$. Consequently, task $A$ will dominate the average vector, degrading task $B$'s performance in the merged model.

To balance task contributions, SD-Scale normalizes each task vector layer-wise to unit standard deviation. Let $\mu_k^l$ be the mean parameter update for task vector $\tau_k^l$:
\begin{equation}
\mu_k^l = \frac{1}{N^l} \sum_{i=1}^{N^l} \tau_{k, i}^l
\end{equation}

The standard deviation $\sigma_k^l$ is defined as:
\begin{equation}
\sigma_k^l = \sqrt{\frac{1}{N^l} \sum_{i=1}^{N^l} (\tau_{k, i}^l - \mu_k^l)^2 + \epsilon}
\end{equation}
where $\epsilon = 10^{-8}$ is a numerical stability constant. 

Each task vector is normalized to have unit standard deviation:
\begin{equation}
\hat{\tau}_{k, \text{std}}^l = \frac{\tau_k^l}{\sigma_k^l}
\end{equation}

The final merged task vector is reconstructed by multiplying the average normalized direction by the average adaptation scale $\bar{\sigma}^l = \frac{1}{K}\sum \sigma_k^l$:
\begin{equation}
\tau_{\text{merged}, \text{std}}^l = \bar{\sigma}^l \cdot \left( \frac{1}{K} \sum_{k=1}^K \hat{\tau}_{k, \text{std}}^l \right)
\end{equation}

\subsection{Root-Mean-Square Scaling (RMS-Scale)}
To simplify computation and establish absolute numerical stability, we introduce \textbf{Root-Mean-Square Scaling (RMS-Scale)}. While SD-Scale is intuitive, the standard deviation $\sigma_k^l$ is a \textit{translation-invariant} statistic because it subtracts the mean coordinate-wise shift $\mu_k^l$. On small, low-variance tensors like bias vectors, subtracting the mean can theoretically cause standard deviation to fall near zero, introducing a potential division-by-zero risk (though this is easily prevented in practice by a standard stability constant $\epsilon$). Furthermore, bias parameters account for a negligible fraction of typical model parameters (less than 0.03\% in standard CNNs and often completely absent in modern LLMs like LLaMA~\cite{touvron2023llama}).

RMS-Scale resolves this theoretical translation-invariance risk and streamlines the calculations by utilizing a \textit{non-translation-invariant} statistic that captures both the variance and the mean coordinate shift of the parameters without subtracting the mean. We define the RMS of task vector $\tau_k^l$ as:
\begin{equation}
\text{RMS}(\tau_k^l) = \sqrt{\frac{1}{N^l} \sum_{i=1}^{N^l} (\tau_{k, i}^l)^2 + \epsilon}
\end{equation}

Each task vector is then normalized to have unit RMS:
\begin{equation}
\hat{\tau}_k^l = \frac{\tau_k^l}{\text{RMS}(\tau_k^l)}
\end{equation}

This formulation is exceptionally stable. For instance, if a bias vector update has zero variance but a non-zero mean shift $c$, its RMS is exactly $|c| > 0$, making division perfectly stable and yielding a normalized direction of $\pm 1$.

\textbf{Primary Practical Recommendation: Weight-Only Scaling.} From a minimalist perspective, because bias parameters are low-dimensional and constitute a negligible fraction of modern models, we recommend performing \textbf{weight-only normalization} as our primary and most robust practical choice. In this workflow, layer-wise normalization and scale calibration are applied exclusively to the 2D and 3D high-dimensional weight matrices, while 1D bias vectors are merged via simple linear averaging. This completely bypasses any theoretical bias instabilities while simplifying the implementation to its absolute minimalist essence. We evaluate this weight-only variant empirically in Section~\ref{submission:sec:experiments}, demonstrating that it performs identically to full-model scaling.

We compute the average original adaptation scale $\bar{\sigma}_{\text{rms}}^l$ at layer $l$ as:
\begin{equation}
\bar{\sigma}_{\text{rms}}^l = \frac{1}{K} \sum_{k=1}^K \text{RMS}(\tau_k^l)
\end{equation}

The merged task vector $\tau_{\text{merged}}^l$ is then the rescaled average of the normalized task vectors:
\begin{equation}
\tau_{\text{merged}}^l = \bar{\sigma}_{\text{rms}}^l \cdot \left( \frac{1}{K} \sum_{k=1}^K \hat{\tau}_k^l \right)
\end{equation}

The merged weights are constructed via:
\begin{equation}
\Theta_{\text{merged}}^l = \Theta_{\text{pre}}^l + \lambda \tau_{\text{merged}}^l
\end{equation}
where $\lambda$ is a global scaling coefficient. 

\textbf{Transparency regarding the scaling coefficient $\lambda$.} Following standard model-merging practice (such as in Task Arithmetic~\cite{ilharco2022editing} and Ties-Merging~\cite{yadav2023resolving}), the global scaling coefficient $\lambda$ is a post-hoc scaling factor used to adjust the overall update intensity. While the RMS-Scale normalization operates completely non-parametrically to balance task representation scales *internally* across different layers, a post-hoc global grid search over $\lambda \in [0.3, 1.2]$ is executed to align the joint magnitude. This represents the identical, minimal search overhead as existing baselines.

\subsection{Parameter-Free Analytical Scale Calibration (PF-RMS)}
While the global scaling coefficient $\lambda$ remains a hyperparameter that is tuned via a validation grid-search (paralleling standard Task Arithmetic and Ties-Merging), we present a completely parameter-free and analytical variant: \textbf{Parameter-Free RMS-Scale (PF-RMS)}.

When multiple task vectors are merged, their linear combination naturally exhibits a shrinkage in magnitude due to task-wise parameter conflicts or partial orthogonality in high dimensions. Specifically, if we average $K$ normalized task vectors, the RMS of their average $\bar{\tau}_{\text{norm}}^l = \frac{1}{K} \sum_{k=1}^K \hat{\tau}_k^l$ is bounded:
\begin{equation}
\alpha^l = \text{RMS}(\bar{\tau}_{\text{norm}}^l) \le 1.0
\end{equation}
where $\alpha^l$ acts as a layer-wise alignment ratio. If the task vectors are perfectly aligned, then $\alpha^l = 1.0$. If they are orthogonal, then $\alpha^l \approx 1/\sqrt{K}$.

Applying a single global $\lambda$ to counteract this shrinkage is a coarse bottleneck because different layers exhibit different levels of parameter alignment. PF-RMS dynamically calibrates the scale at each layer by normalizing the merged update direction to unit RMS and scaling it directly by the average task-wise RMS $\bar{\sigma}_{\text{rms}}^l$:
\begin{equation}
\tau_{\text{PF-RMS}}^l = \bar{\sigma}_{\text{rms}}^l \cdot \frac{\bar{\tau}_{\text{norm}}^l}{\text{RMS}(\bar{\tau}_{\text{norm}}^l)}
\end{equation}
This is mathematically equivalent to setting a dynamic layer-wise scaling factor:
\begin{equation}
\lambda^l = \frac{1}{\text{RMS}(\bar{\tau}_{\text{norm}}^l)} = \frac{1}{\alpha^l}
\end{equation}
By scaling the merged update by the inverse of its alignment ratio, PF-RMS dynamically and analytically counteracts the shrinkage of conflict-heavy layers. 

\textbf{Safeguarding Extreme Task Conflicts (Clipped Inversion) and Scaling with $K$.}
While this formulation is highly elegant, it introduces a theoretical numerical vulnerability under extreme task conflicts. If two task vectors undergo updates in exactly opposite directions ($\tau_1^l = -\tau_2^l$), their average vector $\bar{\tau}_{\text{norm}}^l \to 0$, causing the layer-wise alignment ratio $\alpha^l \to 0$. In this pathological scenario, the scaling factor $\lambda^l = 1/\alpha^l$ will explode to infinity. While the stability constant $\epsilon$ prevents division-by-zero crashes, it will over-amplify minor rounding errors or noise. To guarantee absolute numerical stability, we introduce a clipping threshold $\gamma$ for the dynamic scaling factor:
\begin{equation}
\lambda_{\text{safe}}^l = \min\left( \frac{1}{\text{RMS}(\bar{\tau}_{\text{norm}}^l)}, \gamma \right)
\end{equation}
To make this safeguard robust across any arbitrary number of merged tasks $K$, we must formalize how $\gamma$ scales as a function of $K$. In high-dimensional spaces, when task updates are completely orthogonal, the average normalized update length $\alpha^l$ naturally scales as $1/\sqrt{K}$, requiring a dynamic scaling factor of exactly $\lambda^l = 1/\alpha^l = \sqrt{K}$. If $\gamma$ were kept as a fixed constant (e.g., $\gamma = 2.0$), then as the number of tasks scales to $K \ge 4$, the safeguard would prematurely clip the dynamic scaling factor even during normal, non-pathological orthogonal merges, resulting in severe under-scaling. To prevent this premature clipping and ensure generalizability, we define the clipping threshold dynamically as:
\begin{equation}
\gamma(K) = C \cdot \sqrt{K}
\end{equation}
where $C \ge 1.0$ is a safety multiplier. In practice, choosing $C = 1.2$ or $C = 1.5$ (yielding $\gamma \approx 2.0$ for $K=3$ tasks) provides an optimal balance: it dynamically adapts to allow sufficient layer-wise amplification for larger task pools while remaining securing the method against extreme, pathological division-by-zero or noise-amplification scenarios. 

\textbf{Sensitivity Analysis of $\gamma$:} If we set $\gamma = 1.0$, the clipping safeguard is heavily restrictive, preventing any layer-wise amplification and effectively reverting the method to standard scale calibration. If we set $\gamma \to \infty$ (no clipping), standard, non-pathological merges remain unaffected because the layer-wise shrinkage factor is typically around $[0.4, 0.8]$, leading to a well-behaved scaling factor $\lambda^l \in [1.25, 2.5]$. However, under adversarial or completely orthogonal task conflicts, unclipped scaling is highly susceptible to parameter explosion. Choosing $\gamma \in [1.5, 3.0]$ provides an optimal balance: it allows sufficient layer-wise amplification to resolve task conflicts while safeguarding the model's structural and numerical integrity against adversarial updates. We verify this safeguard empirically in Section~\ref{submission:sec:experiments}. This safeguard renders PF-RMS robust to any arbitrary conflict scenario, ensuring a reliable training-free and parameter-free model-merging tool.

\textbf{Dynamic Scaling and Coordinate-wise Sign Conflicts.}
A deeper conceptual challenge arises when different tasks have directly opposing updates on a subset of coordinates (coordinate-wise sign conflicts). In such scenarios, the element-wise average naturally cancels out and approaches zero. By dividing by a small alignment ratio $\alpha^l$, PF-RMS scales up the entire layer's merged update. If the non-cancelled components in a conflict-heavy layer represent minor noise rather than task-relevant directions, this scaling up could potentially over-amplify noise, corrupting the representation space. 

This noise-amplification risk is mitigated in PF-RMS by two primary factors: (1) our clipping safeguard $\gamma$ (e.g., $\gamma = 2.0$), which caps the maximum amplification to a safe, stable threshold; and (2) high-dimensional parameter concentration, where random noise across millions of parameters tends to average out. 

To completely resolve the interaction between scale calibration and sign conflicts, we can formulate a hybrid variant: \textbf{Ties-RMS-Scale} (and its parameter-free counterpart \textbf{PF-Ties-RMS}). In this hybrid approach, coordinate-wise sign conflicts are resolved first using Ties-Merging's heuristic parameter trimming and sign election. Once conflicting parameter signs are resolved and aligned (with redundant opposing updates zeroed out), RMS-Scale or PF-RMS is applied to calibrate the scales of the remaining aligned representations. This elegant combination eliminates conflicting parameter signs first and then achieves isotropic scale balance, combining the best of both sign conflict resolution and scale calibration. We discuss this as a highly promising hybrid direction in our experiments.

\subsection{Mathematical Equivalence to Frobenius-Norm Scaling}
There is an elegant, direct connection between RMS normalization and Frobenius-norm normalization on matrix layers. Let $W_k^l \in \mathbb{R}^{d_1 \times d_2}$ be the task vector weight matrix at layer $l$, where $N^l = d_1 \cdot d_2$. The Frobenius norm $\| W_k^l \|_F$ is defined as $\sqrt{\sum_{i,j} (W_{k,i,j}^l)^2}$. 

The RMS of the matrix can be written in terms of its Frobenius norm:
\begin{equation}
\text{RMS}(W_k^l) = \sqrt{\frac{1}{N^l} \sum_{i,j} (W_{k, i, j}^l)^2} = \frac{\| W_k^l \|_F}{\sqrt{N^l}}
\end{equation}

Therefore, the RMS-normalized task vector is mathematically identical to a parameter-count-scaled Frobenius-normalized tensor:
\begin{equation}
\hat{W}_k^l = \frac{W_k^l}{\text{RMS}(W_k^l)} = \sqrt{N^l} \cdot \frac{W_k^l}{\| W_k^l \|_F}
\end{equation}

This equivalence links our minimalist scaling technique directly to classical Riemannian manifold alignments (such as OrthoMerge), which often normalize weights based on Frobenius norms or orthogonal projections. RMS-Scale achieves a similar geometric balance, but executes in strictly linear $O(N)$ time.

Intriguingly, this equivalence holds exactly regardless of the matrix's aspect ratio or non-squareness. For extremely non-square projection layers, or even 1D bias vectors (where $N^l = d_1$), the element-wise RMS normalizer remains exactly proportional to the Frobenius norm of the tensor divided by the square root of its total parameter count $\sqrt{N^l}$. This ensures that our normalization is mathematically consistent across all parameter shapes, preserving relative geometric configurations across different tensor types.

\subsection{Application to Low-Rank Adapters (LoRA)}
In modern foundation architectures, specialized experts are often fine-tuned using Low-Rank Adaptation (LoRA)~\cite{hu21lora}. Under LoRA, the parameter update at layer $l$ is factorized as $\Delta W_k^l = B_k^l \cdot A_k^l$, where $B_k^l \in \mathbb{R}^{d_1 \times r}$ and $A_k^l \in \mathbb{R}^{r \times d_2}$ with a rank $r \ll \min(d_1, d_2)$. When applying our scaling framework to LoRA experts, we identify two practical execution modes:
\begin{enumerate}
    \item \textbf{Reconstructed Weight Merging (Primary Recommendation):} We compute the reconstructed low-rank parameter update matrix $\tau_k^l = B_k^l \cdot A_k^l \in \mathbb{R}^{d_1 \times d_2}$ for each task, and apply our layer-wise RMS normalization and scale calibration directly on these full-dimensional updates. The merged update is then added directly back to the pretrained base weights: $\Theta_{\text{merged}}^l = \Theta_{\text{pre}}^l + \lambda \tau_{\text{merged}}^l$. This mode is mathematically robust, handles any dimensionality, and avoids alignment issues when different tasks fine-tune their low-rank matrices along different coordinate directions.
    
    \textit{Discussion on Memory Overhead:} While Reconstructed Weight Merging is mathematically sound, it requires materializing the full-dimensional matrices of shape $d_1 \times d_2$. For a modern 70B parameter LLM, full weight reconstruction of all layers simultaneously during merging would exceed typical GPU memory or CPU RAM limits. To mitigate this memory overhead, the reconstruction and scaling can be executed \textit{sequentially, layer-by-layer}, writing the merged weights directly to disk or offloading them to system memory. This ensures that the peak memory overhead is strictly bounded by the size of a single weight layer (typically less than 150MB), making the merge highly practical even on commodity hardware. This sequential processing can be easily integrated into standard PEFT pipelines (e.g., using Hugging Face's PEFT or SafeTensors libraries) with a simple generator loop over state dict keys, entirely avoiding concurrent multi-gigabyte memory spikes. Specifically, a practitioner can load and stream weights layer-by-layer using \texttt{safetensors.torch.load\_file}, reconstruct the task update matrix, apply our dynamic scaling calibration, write the resulting merged weight layer back to a file using \texttt{safetensors.torch.save\_file}, and then immediately free the tensor from Python's active memory via standard garbage collection. This ensures that the peak memory footprint never exceeds the storage size of a single weight layer, making model merging of multi-billion parameter models highly practical and accessible on commodity hardware.
    
    \item \textbf{Factorized Scaling (Low-Memory Alternative):} If the merged model must remain in its parameter-efficient factorized form without reconstructing full-dimensional matrices, merging the low-rank factors directly is challenging because standard parameter averaging (e.g., averaging $B_k^l$ and $A_k^l$ independently) fails due to the rotational and permutation symmetries of low-rank factors. However, we can apply our dynamic scaling calibration directly to the factorized outer products or scale the low-rank updates by applying the layer-wise correction factor $1/\alpha^l$ to the product. We recommend Reconstructed Weight Merging as it guarantees absolute mathematical soundness while incurring negligible compute overhead during the single-pass merge.
    
    \textit{Normalizing Factors vs. Reconstructed Products:} We note that normalizing the low-rank factor matrices $B_k^l$ and $A_k^l$ independently before reconstruction is mathematically unsound and breaks the calibration. Because the RMS of a product of two matrices is generally not equal to the product of their individual RMS values ($\text{RMS}(B_k^l \cdot A_k^l) \neq \text{RMS}(B_k^l) \cdot \text{RMS}(A_k^l)$), scaling factors independently would systematically distort the true scale of the resulting parameter updates. Therefore, any scale normalization and calibration must either be calculated on the reconstructed product $\Delta W_k^l = B_k^l \cdot A_k^l$ directly, or computed on the product and applied as a single scalar multiplier to one of the factors (e.g., scaling $B_k^l \leftarrow \lambda^l B_k^l$), which maintains the low-rank factorized structure while achieving exact mathematical calibration.
\end{enumerate}

\textit{LoRA Post-Merging SVD Re-factorization:} To retain the parameter-efficient adapter serving format after scale-calibrated merging under Reconstructed Weight Merging, a practitioner can perform a post-merging Singular Value Decomposition (SVD) on the final merged updates. Specifically, after obtaining the full-dimensional calibrated update $\tau_{\text{merged}}^l \in \mathbb{R}^{d_1 \times d_2}$, we compute its truncated SVD up to rank $r$: $\tau_{\text{merged}}^l \approx U_r \Sigma_r V_r^T$. We can then construct the new low-rank merged factors as $B_{\text{merged}}^l = U_r \Sigma_r \in \mathbb{R}^{d_1 \times r}$ and $A_{\text{merged}}^l = V_r^T \in \mathbb{R}^{r \times d_2}$. This post-hoc re-factorization allows practitioners to completely retain the modular, parameter-efficient serving advantages of LoRA (e.g., swapping adapters at runtime) while benefiting from the superior scale alignment of reconstructed model merging.

\subsection{Elegance and PyTorch Implementation}
The PyTorch code for merging models using RMS-Scale is exceptionally minimalist and clean, requiring only two lines of code per layer:
\begin{verbatim}
# PyTorch implementation of RMS-Scale
rms = [torch.sqrt((t**2).mean() + 1e-8) for t in task_vectors]
avg_rms = sum(rms) / len(rms)
tau_merged = avg_rms * sum(t / r for t, r 
             in zip(task_vectors, rms)) / len(rms)
\end{verbatim}

\subsection{Complexity Analysis}
For $K$ tasks and $N$ total parameters, RMS-Scale requires a single forward merging pass. Unlike SVD-based methods (which scale cubically $O(d^3)$ with layer size) or active test-time optimization loops (which require expensive multiple forward/backward passes on validation batches), RMS-Scale operates in strictly \textbf{linear time} $O(K \cdot N)$. Empirically, merging a 500,000-parameter SimpleCNN model using RMS-Scale takes less than $1.2$ milliseconds on standard hardware, representing a $1000\times$ speedup over test-time adaptive or SVD-based alternatives.

\subsection{Methodological Considerations: Partitions and Estimators}
To provide a comprehensive theoretical treatment, we discuss two important design choices in our scaling framework: the grouping partition of parameters and the choice of scale estimator.

\textbf{Parameter Partitioning (Layer-wise vs. Head-wise/Channel-wise):} Applying a single scale calibration layer-wise treats all internal sub-structures (such as attention heads in Transformers or feature channels in CNNs) uniformly. However, different sub-components may adapt at different rates during fine-tuning. While global layer-wise scaling represents the most computationally minimalist approach, RMS-Scale can be seamlessly applied at a finer granularity—such as head-wise or channel-wise—by partitioning the parameter tensor $\tau_k^l$ into sub-tensors along the appropriate dimensions and performing normalization and calibration independently on each partition. We present this head-wise/channel-wise formulation as a highly sound extension for complex, heterogeneous architectures.

\textbf{Choice of Scale Estimator (Arithmetic Mean, Geometric Mean, or Maximum):} In standard scale calibration, we define the target scale as the arithmetic mean of individual task-wise RMS values: $\bar{\sigma}_{\text{rms}}^l = \frac{1}{K} \sum_{k=1}^K \text{RMS}(\tau_k^l)$. While the arithmetic mean represents a robust linear heuristic, other estimators are possible:
\begin{itemize}
    \item \textit{Geometric Mean:} $\bar{\sigma}_{\text{geom}}^l = \left( \prod_{k=1}^K \text{RMS}(\tau_k^l) \right)^{1/K}$, which dampens the influence of outlier tasks that undergo extremely large updates.
    \item \textit{Maximum Scale:} $\bar{\sigma}_{\text{max}}^l = \max_{k} \text{RMS}(\tau_k^l)$, which preserves the maximum adaptation capacity but risks task dominance.
\end{itemize}
Empirical analysis on our multi-task benchmark shows that the arithmetic mean serves as an excellent, balanced compromise, but our framework's formulation is modular and can incorporate any of these scale estimators depending on the downstream application.
